% !TEX program = pdflatex
% THEME 1 — Moyen — Corrigé
\documentclass[11pt,a4paper]{article}
\usepackage[utf8]{inputenc}
\usepackage[T1]{fontenc}
\usepackage{lmodern}
\usepackage[french]{babel}
\usepackage{amsmath,amssymb,mathtools}
\usepackage{icomma}
\usepackage{siunitx}
\sisetup{output-decimal-marker={,},per-mode=symbol}
\DeclareSIUnit{\atm}{atm}
\DeclareSIUnit{\bar}{bar}
\DeclareSIUnit{\mmHg}{mmHg}
\usepackage[version=4]{mhchem}
\usepackage{xcolor}
\usepackage{colortbl}
\usepackage[most]{tcolorbox}
\usepackage{enumitem}
\usepackage{array}
\usepackage{booktabs}
\usepackage{fancyhdr}
\usepackage[a4paper,margin=2.1cm,headheight=26pt,headsep=10pt]{geometry}

\definecolor{primary}{RGB}{150,98,162}
\definecolor{primarydark}{RGB}{92,52,110}
\definecolor{excol}{RGB}{0,135,90}
\definecolor{attncol}{RGB}{200,40,55}

\newtcolorbox{rep}[1][]{enhanced,breakable,colback=excol!5,colframe=excol,
  boxrule=0.8pt,arc=2pt,left=3mm,right=3mm,top=2mm,bottom=2mm,
  fonttitle=\bfseries,coltitle=white,title={#1}}

\pagestyle{fancy}\fancyhf{}
\fancyhead[L]{\small\textcolor{primarydark}{\textbf{Fiche 7} — Loi des gaz}}
\fancyhead[R]{\small\textcolor{primarydark}{\textsc{Thème 1 · Moyen · Corrigé}}}
\fancyfoot[C]{\small\thepage}
\renewcommand{\headrulewidth}{0.8pt}
\renewcommand{\headrule}{\hbox to\headwidth{\color{primary}\leaders\hrule height \headrulewidth\hfill}}
\setlength{\parindent}{0pt}\setlength{\parskip}{3pt}

\newcounter{exo}
\newenvironment{cor}[1][]{\refstepcounter{exo}\medskip
  \noindent{\color{primarydark}\bfseries Corrigé \theexo.}\ \textit{#1}\par\nopagebreak\smallskip}{\par}
\newcommand{\rf}[1]{\textbf{\textcolor{attncol}{#1}}}

\begin{document}

\begin{center}
{\Large\bfseries\color{primarydark} Thème 1 — Niveau Moyen — Corrigé}
\end{center}
\textcolor{primary}{\hrule height 1pt}
\medskip

\begin{cor}[Bombe aérosol rigide]
\textbf{(a)} La bombe est \emph{rigide} : le volume $V$ ne change pas (transformation \textbf{isochore}),
et la quantité de gaz est constante. La loi adaptée est $\dfrac{p_1}{T_1}=\dfrac{p_2}{T_2}$.\\
\textbf{(b)} On convertit d'abord : $T_1=27+273=\qty{300}{\kelvin}$, $T_2=127+273=\qty{400}{\kelvin}$.
\[ p_2=p_1\frac{T_2}{T_1}=\qty{2.0}{\atm}\times\frac{\qty{400}{\kelvin}}{\qty{300}{\kelvin}}
     =\rf{\qty{2.67}{\atm}}. \]
La pression grimpe de $2{,}0$ à $\approx\qty{2.7}{\atm}$ : c'est pourquoi on ne jette jamais une bombe
aérosol au feu (risque d'explosion).
\end{cor}

\begin{cor}[Compression avec changement d'unité]
Température constante $\Rightarrow$ \textbf{Boyle} : $p_1V_1=p_2V_2$. Les volumes étant tous deux en
\si{\milli\litre}, pas besoin de les convertir (l'unité se simplifie).
\[ p_2=\frac{p_1V_1}{V_2}=\frac{\qty{100}{\kilo\pascal}\times\qty{500}{\milli\litre}}{\qty{250}{\milli\litre}}
     =\rf{\qty{200}{\kilo\pascal}}. \]
En atmosphères : $\dfrac{200}{101}\approx\rf{\qty{1.98}{\atm}}\approx\qty{2}{\atm}$. On a divisé $V$ par
$2$, donc doublé $p$.
\end{cor}

\begin{cor}[Le piège des degrés Celsius]
\textbf{(a)} Pression constante $\Rightarrow$ \textbf{Charles}. On convertit :
$T_1=27+273=\qty{300}{\kelvin}$, $T_2=87+273=\qty{360}{\kelvin}$.
\[ V_2=V_1\frac{T_2}{T_1}=\qty{2.0}{\litre}\times\frac{360}{300}=2{,}0\times1{,}2=\rf{\qty{2.4}{\litre}}. \]
\textbf{(b)} L'erreur est d'avoir utilisé les températures \emph{en degrés Celsius} dans le rapport
$T_2/T_1$. Le rapport $87/27=3{,}2$ n'a aucun sens physique : l'échelle Celsius a une origine
arbitraire ($\qty{0}{\degreeCelsius}$ n'est pas « pas de température »). Multiplier par $3{,}2$ un volume
alors que la température n'a réellement augmenté que de $\qty{300}{\kelvin}$ à $\qty{360}{\kelvin}$
(soit $+20\,\%$) est absurde. \rf{Seuls les kelvins} permettent des rapports corrects.
\end{cor}

\begin{cor}[Loi combinée]
Ici $p$, $V$ et $T$ changent tous, mais $n$ est constant : on utilise la \textbf{loi combinée}
$\dfrac{p_1V_1}{T_1}=\dfrac{p_2V_2}{T_2}$. On isole $p_2$ :
\[ p_2=p_1\times\frac{V_1}{V_2}\times\frac{T_2}{T_1}
     =\qty{1.0}{\atm}\times\frac{10{,}0}{5{,}0}\times\frac{450}{300}
     =1{,}0\times2\times1{,}5=\rf{\qty{3.0}{\atm}}. \]
Deux effets se cumulent et augmentent la pression : la compression (volume $\div2$, soit $p\times2$) et
le chauffage (température $\times1{,}5$, soit $p\times1{,}5$), d'où $p\times3$ au total.
\end{cor}

\begin{cor}[Raisonner en facteurs]
\textbf{(a)} De la loi combinée $\dfrac{p_1V_1}{T_1}=\dfrac{p_2V_2}{T_2}$ on tire
\[ \frac{V_2}{V_1}=\frac{p_1}{p_2}\times\frac{T_2}{T_1}. \]
Ici $\dfrac{T_2}{T_1}=3$ (température $\times3$) et $\dfrac{p_2}{p_1}=2$ (pression doublée), donc
$\dfrac{p_1}{p_2}=\dfrac{1}{2}$.\\
\textbf{(b)} $\dfrac{V_2}{V_1}=\dfrac{1}{2}\times3=\rf{1{,}5}$. Le volume est multiplié par $1{,}5$ :
le gaz s'est \rf{dilaté} (l'effet du chauffage l'emporte sur celui de la compression).
\end{cor}

\end{document}
